\begin{proof}[Proof of \cref{obj:thm:interval-frontier}]
\begin{enumerate}
\item Fix \(\alpha\in(0,1/2)\). Choose the universal risk constant \(C_U\) supplied by \cref{obj:lem:upper-risk}, so that
\[
0<C_U<\infty,
\qquad
\mathbb E_{P_O^{\otimes n}}
\left[
\bigl(\widehat\tau_{\mathrm{MM}}-\tau(P)\bigr)^2
\right]
\le C_U r_{n,d,q}
\]
for every \(n,d\ge1\), every \(q\in[1/2,1]\), and every \(P\in\mathcal M(d,q)\).
The combined interval bound in \cref{obj:lem:prior-reduction} supplies a constant \(c_{\mathrm{lower},\alpha}\), depending only on \(\alpha\), satisfying
\[
c_{\mathrm{lower},\alpha}>0,
\qquad
c_{\mathrm{lower},\alpha}\sqrt{r_{n,d,q}}
\le \mathfrak L^*_{n,d,q,\alpha}
\]
throughout the same parameter range.
% lean: CausalSmith.Stat.MarNearcompleteFrontier.prior_reduction; CausalSmith.Stat.MarNearcompleteFrontier.upper_risk

\item Define the comparison constants by
\[
c_\alpha:=c_{\mathrm{lower},\alpha},
\qquad
C_\alpha:=c_{\mathrm{lower},\alpha}
       +2\sqrt{C_U/\alpha}+1.
\]
Since \(\sqrt{C_U/\alpha}\ge0\), these constants satisfy
\[
0<c_\alpha<C_\alpha<\infty.
\]
Now fix arbitrary positive integers \(n,d\) and \(q\in[1/2,1]\). By \cref{obj:def:rate},
\[
r_{n,d,q}=\frac1n+g_{n,d,q}^{\,2}\ge0,
\qquad
\sqrt{r_{n,d,q}}\ge0.
\]
The preceding lower bound therefore gives the required inequality
\[
c_\alpha\sqrt{r_{n,d,q}}
\le \mathfrak L^*_{n,d,q,\alpha}.
\]
% lean: CausalSmith.Stat.MarNearcompleteFrontier.interval_frontier

\item Apply \cref{obj:lem:honest-interval-construction} with the chosen \(C_U\). For the procedure in \cref{obj:def:mm-interval}, it yields
\[
\widehat I_{\mathrm{MM},\alpha}\in\mathcal C_\alpha(n,d,q)
\]
and, for every \(P\in\mathcal M(d,q)\),
\[
\mathbb E_{P_O^{\otimes n}}
\left[
\operatorname{length}\{\widehat I_{\mathrm{MM},\alpha}(O^n)\}
\right]
\le
\min\left\{2,2\sqrt{\frac{C_Ur_{n,d,q}}{\alpha}}\right\}.
\]

The law class is nonempty. Indeed, the paired construction with all latent masses zero concentrates on its filler label. Write \(P_{\mathrm{fill}}\) for this law, specified by
\[
\begin{gathered}
X=0,\qquad S(0)=S(1)=Y(0)=0,\qquad R=1,\\
A\sim\operatorname{Bernoulli}(1/2),\qquad
Y(1)\sim\operatorname{Bernoulli}\!\left(\frac{(1+q)^2}{8q}\right),
\qquad A\perp Y(1),\\
S=S(A),\qquad Y=Y(A),\qquad RY=R\,Y.
\end{gathered}
\]
The Bernoulli parameter belongs to \([0,1]\), since
\[
0<(1+q)^2\le4\le8q.
\]
This law has one occupied baseline label, independent balanced treatment, consistent realized measurements, and complete arrival. Hence, by \cref{obj:def:law-class},
\[
P_{\mathrm{fill}}\in\mathcal M(d,q).
\]
Taking the supremum of the expected-length bound gives
\[
\sup_{P\in\mathcal M(d,q)}
\mathbb E_{P_O^{\otimes n}}
\left[
\operatorname{length}\{\widehat I_{\mathrm{MM},\alpha}(O^n)\}
\right]
\le
\min\left\{2,2\sqrt{\frac{C_Ur_{n,d,q}}{\alpha}}\right\}.
\]
% lean: CausalSmith.Stat.MarNearcompleteFrontier.honest_interval_construction; CausalSmith.Stat.MarNearcompleteFrontier.pairedClassLaw

\item Because \(C_U/\alpha\ge0\) and \(r_{n,d,q}\ge0\), square-root multiplicativity gives
\[
\frac{C_Ur_{n,d,q}}{\alpha}
=\frac{C_U}{\alpha}\,r_{n,d,q},
\qquad
\sqrt{\frac{C_Ur_{n,d,q}}{\alpha}}
=\sqrt{\frac{C_U}{\alpha}}\sqrt{r_{n,d,q}}.
\]
Consequently,
\[
\min\left\{2,2\sqrt{\frac{C_Ur_{n,d,q}}{\alpha}}\right\}
\le
2\sqrt{\frac{C_U}{\alpha}}\sqrt{r_{n,d,q}}.
\]
The definition of \(C_\alpha\) and positivity of \(c_{\mathrm{lower},\alpha}\) imply
\[
2\sqrt{\frac{C_U}{\alpha}}\le C_\alpha.
\]
Multiplying by the nonnegative quantity \(\sqrt{r_{n,d,q}}\) yields
\[
2\sqrt{\frac{C_U}{\alpha}}\sqrt{r_{n,d,q}}
\le C_\alpha\sqrt{r_{n,d,q}}.
\]
% lean: CausalSmith.Stat.MarNearcompleteFrontier.interval_frontier

\item For every \(I\in\mathcal C_\alpha(n,d,q)\), the endpoint ordering in \cref{obj:def:interval-class} gives
\[
\operatorname{length}\{I(o)\}\ge0
\quad\text{for every }o\in O^n.
\]
If the law-indexed expected lengths of this candidate interval are bounded above, then
\[
0\le
\mathbb E_{P_{\mathrm{fill},O}^{\otimes n}}
\left[\operatorname{length}\{I(O^n)\}\right]
\le
\sup_{P\in\mathcal M(d,q)}
\mathbb E_{P_O^{\otimes n}}
\left[\operatorname{length}\{I(O^n)\}\right].
\]
For an unbounded-above collection, the real-valued supremum is assigned the value zero. Thus, in the unbounded case,
\[
0=
\sup_{P\in\mathcal M(d,q)}
\mathbb E_{P_O^{\otimes n}}
\left[\operatorname{length}\{I(O^n)\}\right].
\]
In either case, the worst-law expected length is nonnegative. The collection of worst-law expected lengths is therefore bounded below by zero and contains the value attained by \(\widehat I_{\mathrm{MM},\alpha}\). Using the infimum in \cref{obj:def:length-risk}, followed by the preceding bounds, gives
\[
\begin{aligned}
\mathfrak L^*_{n,d,q,\alpha}
&\le
\sup_{P\in\mathcal M(d,q)}
\mathbb E_{P_O^{\otimes n}}
\left[
\operatorname{length}\{\widehat I_{\mathrm{MM},\alpha}(O^n)\}
\right]\\
&\le
\min\left\{2,2\sqrt{\frac{C_Ur_{n,d,q}}{\alpha}}\right\}\\
&\le
2\sqrt{\frac{C_U}{\alpha}}\sqrt{r_{n,d,q}}\\
&\le C_\alpha\sqrt{r_{n,d,q}}.
\end{aligned}
\]
Together with the lower bound, this proves the assertion uniformly over the stated parameter range.
% lean: CausalSmith.Stat.MarNearcompleteFrontier.interval_frontier
\end{enumerate}
\end{proof}
